\documentclass[12pt]{amsart}
\pagestyle{plain}

\usepackage{amsmath,amssymb}
\usepackage{enumerate}
\usepackage[shortlabels]{enumitem}
\setlist[enumerate]{itemsep=6pt, topsep=6pt, parsep=0pt}
\usepackage{mathtools}
\usepackage{xcolor}
\usepackage[left=0.75in,right=1in,bottom=0.9in]{geometry}
\usepackage[T1]{fontenc}
\usepackage[parfill]{parskip}
\renewcommand{\baselinestretch}{1}

% SECTION DIVIDERS
\newcommand{\divider}{
  \vspace{-1em}
  \noindent
  \raisebox{-0.15ex}{\textbullet}%
  \hspace{0.5em}%
  \rule[0.5ex]{0.95\textwidth}{1.0pt}%
  \hspace{0.5em}%
  \raisebox{-0.15ex}{\textbullet}%
  \vspace{0em}
}

\title{ESE 2030 QUIZZAM 2 : SOLUTIONS GUIDE}
\author{prof-g}

\begin{document}
\maketitle

\begin{center}
{\bf NOTE TO STUDENTS}
\end{center}

This solutions guide is organized by course week and topic.
Since there are five versions of this quiz with permuted problem orders \emph{and}
permuted answer choices, the guide does \emph{not} identify correct answers by
letter~(A)--(E).
Instead, the correct answer is \textcolor{blue}{\textbf{highlighted in blue}}
directly within each list of choices.
Match the problem text to the one you received, then read the highlighted
choice, the explanation, and the analysis of each distractor.

Each solution includes:
\begin{itemize}
\item The correct answer (highlighted in blue in the choice list)
\item Detailed mathematical explanation
\item Analysis of every wrong choice, including partial credit and common misconceptions (``trick answers'')
\end{itemize}

\newpage

%%%%%%%%%%%%%%%%%%%%%%%%%%%%%%%%%%%%%%%%%%%%%%%%%%%%%%%%%%%
\section*{WEEK 4: BASES AND COORDINATES}
%%%%%%%%%%%%%%%%%%%%%%%%%%%%%%%%%%%%%%%%%%%%%%%%%%%%%%%%%%%

%----------------------------------------------------------
\subsection*{Problem: Similar matrices (same transformation, different bases)}
%----------------------------------------------------------

Two engineers model the same linear transformation $T:\mathbb{R}^n\to\mathbb{R}^n$
using different coordinate systems.
Engineer~1 uses basis~$\mathcal{A}$ and obtains matrix~$A$; Engineer~2 uses
basis~$\mathcal{B}$ and obtains matrix~$B$. Which statement is correct?

\begin{enumerate}[(A)]
\item $A$ and $B$ have the same kernel
\item $A$ and $B$ have the same image
\item $A=B$ if the bases $\mathcal{A}$ and $\mathcal{B}$ are both orthonormal
\item \textcolor{blue}{\textbf{$A$ and $B$ are similar}}
\item $A$ and $B$ are invertible
\end{enumerate}

\textbf{Explanation:}
Since both matrices represent the same transformation in different bases,
there is a change-of-basis matrix $P$ (from $\mathcal{A}$ to $\mathcal{B}$) such that
$B = P^{-1}AP$.
This is the definition of similarity.

\textbf{Partial credit:}
The choices stating that $A$ and $B$ have the same kernel (or the same image)
reflect a correct intuition --- the abstract kernel and image of $T$ are
basis-independent --- but confuse the abstract subspace of $\mathbb{R}^n$ with its
coordinate representation.
$\ker(A)$ and $\ker(B)$ are sets of coordinate vectors in \emph{different}
coordinate systems, so they are generically different subsets of $\mathbb{R}^n$.

\textbf{Trick answer:}
Having both bases be orthonormal plays no special role; two orthonormal bases
still generically yield different matrices.

\textbf{Trick answer:}
$T$ need not be invertible (similarity preserves rank but not invertibility),
so neither $A$ nor $B$ need be invertible.

\divider

%----------------------------------------------------------
\subsection*{Problem: Coordinate isomorphism and linear dependence}
%----------------------------------------------------------

Let $V$ be a 3-dimensional vector space with basis
$\mathcal{B} = \{\mathbf{b}_1, \mathbf{b}_2, \mathbf{b}_3\}$.
Three vectors in $V$ have coordinates
$$[\mathbf{u}]_\mathcal{B} = \begin{pmatrix} 1 \\ 2 \\ -1 \end{pmatrix}, \qquad
  [\mathbf{v}]_\mathcal{B} = \begin{pmatrix} 3 \\ 0 \\ 1 \end{pmatrix}, \qquad
  [\mathbf{w}]_\mathcal{B} = \begin{pmatrix} 4 \\ 2 \\ 0 \end{pmatrix}.$$
Which statement about $\{\mathbf{u},\mathbf{v},\mathbf{w}\}$ is correct?

\begin{enumerate}[(A)]
\item They are linearly independent because $V$ is 3-dimensional and we have exactly 3 vectors.
\item Whether they are linearly dependent or independent cannot be determined without knowing $\mathcal{B}$ explicitly.
\item They are linearly independent because $\mathcal{B}$ is a basis.
\item \textcolor{blue}{\textbf{They are linearly dependent because $[\mathbf{w}]_\mathcal{B} = [\mathbf{u}]_\mathcal{B} + [\mathbf{v}]_\mathcal{B}$.}}
\item None of the above.
\end{enumerate}

\textbf{Explanation:}
The coordinate map $\mathbf{x}\mapsto[\mathbf{x}]_\mathcal{B}$ is a
\emph{linear isomorphism} from $V$ to $\mathbb{R}^3$.
Because it is linear, every linear relation among vectors in $V$ is faithfully
reflected in their coordinate vectors.
Since $(1,2,-1)^T + (3,0,1)^T = (4,2,0)^T$, we have $\mathbf{w} = \mathbf{u}+\mathbf{v}$ in $V$.
No explicit knowledge of $\mathcal{B}$ is needed beyond the fact that it is a basis.

\textbf{Trick answer:}
Having $n$ vectors in an $n$-dimensional space is necessary for a basis but
not sufficient for independence; three coplanar vectors in $\mathbb{R}^3$ can be
dependent.
This conflates a necessary condition with a sufficient one.

\textbf{Partial credit:}
The choice saying independence cannot be determined without knowing $\mathcal{B}$
explicitly shows the right instinct (caution about abstract spaces) but has
not internalized the key theorem: the coordinate isomorphism means that linear
independence can always be settled from the coordinates alone.

\textbf{Trick answer:}
The choice asserting independence because $\mathcal{B}$ is a basis likely reduces
``linear independence'' to pairwise non-proportionality.
The three coordinate vectors are pairwise non-proportional, yet they satisfy
$[\mathbf{u}]_\mathcal{B} + [\mathbf{v}]_\mathcal{B} - [\mathbf{w}]_\mathcal{B} = \mathbf{0}$.

\divider

%----------------------------------------------------------
\subsection*{Problem: Matrix representation and the images of basis vectors}
%----------------------------------------------------------

Let $T: V \to V$ be a linear transformation on an $n$-dimensional vector space
($n \geq 3$), and let $\mathcal{B} = \{\mathbf{b}_1, \ldots, \mathbf{b}_n\}$ be a basis for $V$.
The matrix representation $[T]_\mathcal{B}$ has its \emph{third column} equal to the zero vector.
Which conclusion follows?

\begin{enumerate}[(A)]
\item $\operatorname{coker}(T)$ contains $\mathbf{b}_3$
\item $\operatorname{im}(T)$ does not contain $\mathbf{b}_3$
\item \textcolor{blue}{\textbf{$\ker(T)$ contains $\mathbf{b}_3$}}
\item $\operatorname{rank}(T) = n - 1$
\item None of the above must be true
\end{enumerate}

\textbf{Explanation:}
The $j$-th column of $[T]_\mathcal{B}$ records $[T(\mathbf{b}_j)]_\mathcal{B}$.
If the third column is $\mathbf{0}$, then $[T(\mathbf{b}_3)]_\mathcal{B} = \mathbf{0}$,
which means $T(\mathbf{b}_3) = \mathbf{0}$.
Therefore $\mathbf{b}_3 \in \ker(T)$, and $T$ is not injective.

\textbf{Trick answer:}
The choice claiming $\operatorname{im}(T)$ does not contain $\mathbf{b}_3$ confuses the
column space of the coordinate matrix with the image of $T$ in $V$.
The other columns of $[T]_\mathcal{B}$ can have nonzero third components, meaning
some $T(\mathbf{b}_j)$ for $j\neq 3$ may involve $\mathbf{b}_3$.
For example, if the first column of $[T]_\mathcal{B}$ is $(0,0,1,0,\ldots)^T$, then
$T(\mathbf{b}_1) = \mathbf{b}_3 \in \operatorname{im}(T)$.

\textbf{Partial credit:}
The choice $\operatorname{rank}(T) = n-1$ correctly recognizes that rank is reduced,
but overstates: one zero column gives $\operatorname{rank}\leq n-1$, not necessarily
$=n-1$, since the remaining $n-1$ columns could also be linearly dependent.

\divider

%----------------------------------------------------------
\subsection*{Problem: Coordinates in the Lagrange basis}
%----------------------------------------------------------

Consider $\mathcal{P}_2$, with the Lagrange basis $\mathcal{L} = \{L_0,L_1,L_2\}$
associated to the nodes $x=-1,0,1$:
each $L_i$ equals $1$ at its own node and $0$ at the others.
$$L_0(-1)=1,\; L_0(0)=0,\; L_0(1)=0; \quad
  L_1(-1)=0,\; L_1(0)=1,\; L_1(1)=0; \quad
  L_2(-1)=0,\; L_2(0)=0,\; L_2(1)=1.$$
What is $[p]_\mathcal{L}$ for $p(x) = x^2 - x + 2$?

\begin{enumerate}[(A)]
\item $\begin{pmatrix}2\\-1\\1\end{pmatrix}$
\item $\begin{pmatrix}2\\2\\4\end{pmatrix}$
\item $\begin{pmatrix}1\\-1\\2\end{pmatrix}$
\item \textcolor{blue}{$\boldsymbol{\begin{pmatrix}4\\2\\2\end{pmatrix}}$}
\item Cannot be determined without computing $L_0$, $L_1$, $L_2$ explicitly
\end{enumerate}

\textbf{Explanation:}
Write $p = c_0 L_0 + c_1 L_1 + c_2 L_2$.
Evaluating at node $x=-1$: the defining property gives $p(-1) = c_0$;
at $x=0$: $p(0)=c_1$; at $x=1$: $p(1) = c_2$.
So the coordinates are simply the evaluations of $p$ at the nodes:
$$p(-1) = 1+1+2 = 4, \qquad p(0) = 0-0+2=2, \qquad p(1)=1-1+2=2.$$
Thus $[p]_\mathcal{L} = (4,2,2)^T$.
The elegant key insight: the Lagrange basis is designed so that
\emph{coordinates are evaluations at the nodes} --- no explicit formula for the
$L_i$ is needed.

\textbf{Trick answer:}
The choice $(2,-1,1)^T$ reads off the coefficients from the standard-basis
representation $p(x) = 2\cdot 1 + (-1)\cdot x + 1\cdot x^2$, conflating the
monomial basis $\{1,x,x^2\}$ with $\mathcal{L}$.
This is the most common error: assuming ``coordinates'' always means
coefficients in the standard form.

\textbf{Partial credit:}
The choice $(2,2,4)^T$ understands the evaluation idea correctly (coordinates
$=$ evaluations at the nodes) but evaluates at the wrong node set $\{0,1,2\}$
instead of $\{-1,0,1\}$.
The conceptual understanding is sound; the error is in reading the problem
parameters.

\textbf{Trick answer:}
The choice $(1,-1,2)^T$ reads the coefficients in descending-degree order from
$x^2-x+2$: a different flavor of the monomial-basis confusion.

\textbf{Trick answer:}
Saying it ``cannot be determined without computing $L_0,L_1,L_2$ explicitly''
misses the defining property of the Lagrange basis.
A student who picks this has not seen that the definition directly determines
the coordinates without any system of equations.

\divider

%----------------------------------------------------------
\subsection*{Problem: Basis dependence of coordinate representation}
%----------------------------------------------------------

Let $\mathcal{B}$ and $\mathcal{C}$ be two \emph{different} ordered bases for $\mathbb{R}^3$.
Suppose vectors $\mathbf{v}$ and $\mathbf{w}$ satisfy
$$[\mathbf{v}]_{\mathcal{B}} = [\mathbf{w}]_{\mathcal{C}} = \begin{pmatrix} 2 \\ -1 \\ 3 \end{pmatrix}.$$
Which statement is correct?

\begin{enumerate}[(A)]
\item $\mathbf{v} = \mathbf{w}$, since their coordinate representations are identical.
\item \textcolor{blue}{\textbf{$\mathbf{v}$ and $\mathbf{w}$ need not be equal.}}
\item $\|\mathbf{v}\| = \|\mathbf{w}\|$, since the norm is determined by the coordinate entries.
\item $\mathbf{v} = \mathbf{w}$ if and only if both $\mathcal{B}$ and $\mathcal{C}$ are orthonormal.
\item The equation $[\mathbf{v}]_{\mathcal{B}} = [\mathbf{w}]_{\mathcal{C}}$ forces $\mathcal{B} = \mathcal{C}$.
\end{enumerate}

\textbf{Explanation:}
Writing $\mathcal{B} = \{\mathbf{b}_1,\mathbf{b}_2,\mathbf{b}_3\}$ and
$\mathcal{C}=\{\mathbf{c}_1,\mathbf{c}_2,\mathbf{c}_3\}$, we have
$\mathbf{v} = 2\mathbf{b}_1 - \mathbf{b}_2 + 3\mathbf{b}_3$ and
$\mathbf{w} = 2\mathbf{c}_1 - \mathbf{c}_2 + 3\mathbf{c}_3$.
Since $\mathcal{B}\neq\mathcal{C}$, these are generically different vectors.
The coordinate vector is not the vector itself; it is a representation that
depends on the chosen basis.

\textbf{Trick answer:}
Claiming $\mathbf{v}=\mathbf{w}$ is the most common misconception about
coordinates: conflating a vector with its coordinate representation.
``Same numbers'' does not mean ``same vector'' when the bases differ.

\textbf{Trick answer:}
The norm of a vector equals the norm of its coordinate vector \emph{only when
the basis is orthonormal}.
For a general basis, $\|\mathbf{v}\|^2$ involves cross terms from the Gram matrix
of $\mathcal{B}$, so identical coordinate entries in different bases need not give
equal norms.

\textbf{Trick answer:}
Orthonormality is not the relevant condition for equality.
$\mathbf{v} = \mathbf{w}$ if and only if $2\mathbf{b}_1 - \mathbf{b}_2 + 3\mathbf{b}_3
= 2\mathbf{c}_1 - \mathbf{c}_2 + 3\mathbf{c}_3$, which depends on the specific basis
vectors, not on whether they are orthonormal.

\divider

%----------------------------------------------------------
\subsection*{Problem: The change-of-basis matrix}
%----------------------------------------------------------

Let $\mathcal{B}$ and $\mathcal{C}$ be two bases for $\mathbb{R}^n$, and let $P$ be the matrix
satisfying $[\mathbf{v}]_\mathcal{C} = P\,[\mathbf{v}]_\mathcal{B}$ for all $\mathbf{v} \in \mathbb{R}^n$.
What are the \emph{columns} of $P$?

\begin{enumerate}[(A)]
\item The vectors in $\mathcal{C}$, expressed in standard coordinates.
\item The vectors in $\mathcal{B}$, expressed in standard coordinates.
\item \textcolor{blue}{\textbf{The vectors in $\mathcal{B}$, expressed in $\mathcal{C}$-coordinates.}}
\item The vectors in $\mathcal{C}$, expressed in $\mathcal{B}$-coordinates.
\item The standard basis vectors, expressed in $\mathcal{C}$-coordinates.
\end{enumerate}

\textbf{Explanation:}
Plug in $\mathbf{v} = \mathbf{b}_j$ (the $j$-th vector of $\mathcal{B}$).
Its $\mathcal{B}$-coordinates are $\mathbf{e}_j$, the $j$-th standard basis vector.
Therefore $P\mathbf{e}_j = [\mathbf{b}_j]_\mathcal{C}$: the $j$-th column of $P$ is the
$j$-th basis vector of $\mathcal{B}$, expressed in $\mathcal{C}$-coordinates.

\textbf{Trick answer:}
The choice giving the vectors of $\mathcal{C}$ in $\mathcal{B}$-coordinates describes
the columns of $P^{-1}$, not $P$.
Students who get the direction of conversion backwards pick this.

\textbf{Trick answer:}
The choice giving the vectors of $\mathcal{B}$ in standard coordinates describes the
matrix that converts from $\mathcal{B}$-coordinates to standard coordinates, not to
$\mathcal{C}$-coordinates.

\textbf{Trick answer:}
Similarly, the vectors of $\mathcal{C}$ in standard coordinates convert from
$\mathcal{C}$-coordinates to standard, the reverse direction again.

\divider

%%%%%%%%%%%%%%%%%%%%%%%%%%%%%%%%%%%%%%%%%%%%%%%%%%%%%%%%%%%
\section*{WEEK 5: INNER PRODUCT SPACES}
%%%%%%%%%%%%%%%%%%%%%%%%%%%%%%%%%%%%%%%%%%%%%%%%%%%%%%%%%%%

%----------------------------------------------------------
\subsection*{Problem: Axioms for a valid inner product}
%----------------------------------------------------------

Which of the following defines a valid inner product on $\mathbb{R}^2$?

\begin{enumerate}[(A)]
\item $\langle \mathbf{u}, \mathbf{v} \rangle = u_1 v_1 - u_2 v_2$
\item $\langle \mathbf{u}, \mathbf{v} \rangle = u_1 v_2 + u_2 v_1$
\item \textcolor{blue}{$\boldsymbol{\langle \mathbf{u}, \mathbf{v} \rangle = 2u_1 v_1 + u_1 v_2 + u_2 v_1 + 3u_2 v_2}$}
\item $\langle \mathbf{u}, \mathbf{v} \rangle = (u_1 + u_2)(v_1 + v_2)$
\item $\langle \mathbf{u}, \mathbf{v} \rangle = u_1 v_1 + u_2 v_2 + 1$
\end{enumerate}

\textbf{Explanation:}
The correct choice corresponds to $\langle\mathbf{u},\mathbf{v}\rangle = \mathbf{u}^T M\mathbf{v}$
with $M = \begin{bmatrix}2&1\\1&3\end{bmatrix}$.
This matrix is symmetric (so the symmetry axiom holds), and it is positive
definite: $\det(M) = 5 > 0$ and $\operatorname{tr}(M) = 5 > 0$ give positive
eigenvalues.
All three inner product axioms are satisfied.

\textbf{Trick answer:}
The choice $u_1v_1 - u_2v_2$ is symmetric and bilinear, but fails positive
definiteness: evaluating at $(0,1)$ gives $-1 < 0$.

\textbf{Trick answer:}
The choice $u_1v_2 + u_2v_1$ is bilinear and symmetric, but fails positive
definiteness: evaluating at $(1,0)$ gives $0$ for a nonzero vector.

\textbf{Trick answer:}
The choice $(u_1+u_2)(v_1+v_2)$ is symmetric and bilinear, but fails positive
definiteness: $(1,-1)$ is a nonzero vector with $\langle(1,-1),(1,-1)\rangle = 0$.
This is a rank-1 form.

\textbf{Wrong:}
The choice $u_1v_1 + u_2v_2 + 1$ fails $\langle\mathbf{0},\mathbf{0}\rangle = 0$ (it gives $1$)
and also fails linearity.

\divider

%----------------------------------------------------------
\subsection*{Problem: Weighted inner products and orthogonality}
%----------------------------------------------------------

In $\mathbb{R}^2$, consider the weighted inner product
$\langle \mathbf{u}, \mathbf{v} \rangle_M = \mathbf{u}^T M \mathbf{v}$ with
$M = \begin{bmatrix} 3 & 1 \\ 1 & 4 \end{bmatrix}$.
Let $\mathbf{a} = (1,-2)^T$ and $\mathbf{b} = (2,1)^T$.
Which statement is correct?

\begin{enumerate}[(A)]
\item $\mathbf{a}$ and $\mathbf{b}$ are orthogonal under both the standard dot product and $\langle\cdot,\cdot\rangle_M$.
\item $\mathbf{a}$ and $\mathbf{b}$ are orthogonal under $\langle\cdot,\cdot\rangle_M$ but not under the standard dot product.
\item \textcolor{blue}{\textbf{$\mathbf{a}$ and $\mathbf{b}$ are orthogonal under the standard dot product but not under $\langle\cdot,\cdot\rangle_M$.}}
\item $\mathbf{a}$ and $\mathbf{b}$ are not orthogonal under either inner product, but $\|\mathbf{a}\|_M = \|\mathbf{a}\|$.
\item None of the above.
\end{enumerate}

\textbf{Explanation:}
Standard dot product: $\mathbf{a}\cdot\mathbf{b} = (1)(2)+(-2)(1) = 0$. So $\mathbf{a}\perp\mathbf{b}$
under the standard inner product.

Weighted inner product:
$\langle\mathbf{a},\mathbf{b}\rangle_M = \mathbf{a}^T M\mathbf{b}$.
We compute $M\mathbf{b} = (7,6)^T$, so $\langle\mathbf{a},\mathbf{b}\rangle_M = (1)(7)+(-2)(6) = 7-12=-5\neq 0$.
Therefore $\mathbf{a}$ and $\mathbf{b}$ are \emph{not} $M$-orthogonal.
The core lesson: orthogonality depends on the choice of inner product.
Two vectors perpendicular in the standard geometry need not be perpendicular
when distances are weighted differently.

\textbf{Trick answer:}
Assuming orthogonality under one inner product implies orthogonality under all
inner products is wrong; this is exactly the misconception the problem tests.

\textbf{Trick answer:}
The choice getting the conclusion backwards --- orthogonal under the weighted
product but not the standard --- reverses the two computations.

\textbf{Wrong:}
The choice claiming they are not orthogonal under either product contradicts the
calculation.
Also, $\|\mathbf{a}\| = \sqrt{5}$ while $\|\mathbf{a}\|_M = \sqrt{\mathbf{a}^TM\mathbf{a}} = \sqrt{10}$,
so the norms differ.

\divider

%----------------------------------------------------------
\subsection*{Problem: The Frobenius inner product}
%----------------------------------------------------------

Consider the vector space $M_{m\times n}(\mathbb{R})$ equipped with the Frobenius
inner product $\langle A, B \rangle_F = \operatorname{tr}(A^T B)$.
Which of the following is TRUE?

\begin{enumerate}[(A)]
\item For square matrices, $\|A\|_F = \sqrt{\operatorname{tr}(A^2)}$.
\item $\langle A, B \rangle_F = 0$ if and only if $A^T B = \mathbf{0}$.
\item \textcolor{blue}{$\boldsymbol{\langle A, B \rangle_F = \sum_{i,j} a_{ij}\,b_{ij}}$}
\item This satisfies the definition of an inner product only for square matrices.
\item None of the above.
\end{enumerate}

\textbf{Explanation:}
Expanding the trace:
$\operatorname{tr}(A^T B) = \sum_{j=1}^{n}\sum_{i=1}^{m}a_{ij}\,b_{ij}$,
the standard dot product applied to the $mn$ entries of $A$ and $B$ laid out
as vectors.
This is the key insight: the Frobenius inner product ``vectorizes'' a matrix
and takes the ordinary dot product, so all the machinery of inner product
spaces applies to matrices.

\textbf{Trick answer:}
The first choice confuses $\operatorname{tr}(A)$ with $\|A\|_F$.
The trace sums only the diagonal entries $a_{ii}$, while
$\|A\|_F^2 = \sum_{i,j}a_{ij}^2$ involves \emph{all} entries.
For example, $A = \begin{psmallmatrix}0&3\\0&0\end{psmallmatrix}$
has $\operatorname{tr}(A^2) = 0$ but $\|A\|_F^2 = 9$.

\textbf{Trick answer:}
The condition $A^TB = \mathbf{0}$ (matrix zero) is far stronger than
$\operatorname{tr}(A^TB) = 0$ (a single number being zero).
For example, $I_2$ and the rotation $\begin{psmallmatrix}0&1\\-1&0\end{psmallmatrix}$
are Frobenius-orthogonal (trace of their product is $0$), yet their product is
not the zero matrix.

\textbf{Wrong:}
The Frobenius inner product requires only that $A$ and $B$ have the same
dimensions, not that they be square.

\divider

%----------------------------------------------------------
\subsection*{Problem: Orthogonal matrices --- what is \emph{not} guaranteed}
%----------------------------------------------------------

Let $Q$ be an $n\times n$ orthogonal matrix.
Which of the following is \emph{not} guaranteed to be true?

\begin{enumerate}[(A)]
\item $\|Q\mathbf{x}\| = \|\mathbf{x}\|$ for all $\mathbf{x} \in \mathbb{R}^n$
\item \textcolor{blue}{\textbf{$\det(Q) = 1$}}
\item The rows of $Q$ form an orthonormal set
\item $Q^{-1} = Q^T$
\item $\langle Q\mathbf{x}, Q\mathbf{y} \rangle = \langle \mathbf{x}, \mathbf{y} \rangle$ for all $\mathbf{x}, \mathbf{y}$
\end{enumerate}

\textbf{Explanation:}
From $Q^TQ = I$ we get $(\det Q)^2 = \det(Q^T)\det(Q) = 1$, so $\det(Q) = \pm 1$.
Reflections are orthogonal and have determinant $-1$, so $\det(Q) = 1$ is
not guaranteed.
All other listed properties are guaranteed by the definition $Q^TQ = I$.

\textbf{Wrong:}
Norm preservation follows from
$\|Q\mathbf{x}\|^2 = \langle Q\mathbf{x},Q\mathbf{x}\rangle = \langle\mathbf{x},\mathbf{x}\rangle = \|\mathbf{x}\|^2$.

\textbf{Partial credit:}
Students often memorize that the \emph{columns} are orthonormal and are
uncertain whether the rows are too.
In fact, $Q^TQ = I$ implies $QQ^T = I$ (for square $Q$), which means the rows
are orthonormal as well.
A student who flags row orthonormality as uncertain shows partial understanding.

\textbf{Wrong:}
$Q^{-1} = Q^T$ is the defining equivalent condition.
Inner product preservation follows immediately from it.

\divider

%----------------------------------------------------------
\subsection*{Problem: Orthogonal complements}
%----------------------------------------------------------

Consider $W < \mathbb{R}^7$ with $\dim W = 3$.
Which statement about $W^\perp$ is true?

\begin{enumerate}[(A)]
\item \textcolor{blue}{\textbf{$\mathbb{R}^7 = W \oplus W^\perp$}}
\item $W \cap W^\perp = W$
\item $\dim(W^\perp) = 3$
\item $\dim\left((W^\perp)^\perp\right) = 4$
\item $W$ and $W^\perp$ do not intersect
\end{enumerate}

\textbf{Explanation:}
The orthogonal decomposition theorem guarantees $\mathbb{R}^7 = W \oplus W^\perp$:
every vector decomposes uniquely into a component in $W$ and a component in
$W^\perp$, and these components are orthogonal.
Since $\dim(W) = 3$, we have $\dim(W^\perp) = 7-3=4$.

\textbf{Wrong:}
$W \cap W^\perp = \{\mathbf{0}\}$: a nonzero vector cannot be orthogonal to itself
since $\langle\mathbf{v},\mathbf{v}\rangle > 0$ for $\mathbf{v}\neq\mathbf{0}$.

\textbf{Partial credit:}
The choice $\dim(W^\perp) = 3$ correctly engages with computing $\dim(W^\perp)$
but applies the wrong formula, assuming $\dim(W^\perp) = \dim(W)$ rather than
$n - \dim(W)$.
It shows awareness that the complement has a computable dimension.

\textbf{Trick answer:}
The double complement satisfies $(W^\perp)^\perp = W$, so its dimension is $3$,
not $4$.
Students who think taking $\perp$ twice yields something new pick this.

\textbf{Trick answer:}
$W\cap W^\perp = \{\mathbf{0}\}$, so they \emph{do} intersect (at the zero vector);
saying they ``do not intersect'' is false.

\divider

%----------------------------------------------------------
\subsection*{Problem: Definition of the adjoint operator}
%----------------------------------------------------------

Let $V$ and $W$ be finite-dimensional inner product spaces with inner products
$\langle\cdot,\cdot\rangle_V$ and $\langle\cdot,\cdot\rangle_W$, and let
$T: V \to W$ be linear.
Which of the following correctly defines the adjoint $T^*$?

\begin{enumerate}[(A)]
\item \textcolor{blue}{\textbf{For each $\mathbf{w} \in W$, the vector $T^*(\mathbf{w})$ is the unique element of $V$ satisfying $\langle T(\mathbf{v}), \mathbf{w} \rangle_W = \langle \mathbf{v}, T^*(\mathbf{w}) \rangle_V$ for every $\mathbf{v} \in V$.}}
\item $T^*(\mathbf{w})$ is the vector in $V$ that minimizes $\|T(\mathbf{v}) - \mathbf{w}\|_W$.
\item $T^*$ is the unique linear map $W \to V$ satisfying $T^* \circ T = I_V$.
\item $T^*(\mathbf{w})$ is the orthogonal projection of $\mathbf{w}$ onto $\operatorname{im}(T)$.
\item $T^* = T^{-1}$ when $T$ is invertible, and $T^* = (A^TA)^{-1}A^T$ otherwise.
\end{enumerate}

\textbf{Explanation:}
The adjoint is defined by ``moving $T$ across the inner product'':
$\langle T(\mathbf{v}),\mathbf{w}\rangle_W = \langle\mathbf{v},T^*(\mathbf{w})\rangle_V$
for all $\mathbf{v}\in V$, $\mathbf{w}\in W$.
The Riesz representation theorem guarantees that for each fixed $\mathbf{w}$, the
functional $\mathbf{v}\mapsto\langle T(\mathbf{v}),\mathbf{w}\rangle_W$ is represented by a
unique vector in $V$; that vector is $T^*(\mathbf{w})$.

\textbf{Trick answer:}
The description of minimizing $\|T(\mathbf{v})-\mathbf{w}\|_W$ describes the
least-squares/pseudoinverse solution $T^\dagger\mathbf{w}$, not the adjoint.

\textbf{Trick answer:}
$T^*\circ T = I_V$ would make $T$ an isometry and fails whenever $T$ has a
nontrivial kernel.
This confuses the adjoint with a left inverse.

\textbf{Trick answer:}
The orthogonal projection onto $\operatorname{im}(T)$ is $TT^\dagger$, which maps
$W\to W$, not $W\to V$, so the types do not match.

\textbf{Trick answer:}
The last choice conflates the adjoint with the pseudoinverse.
$T^*$ always exists and equals the transpose $A^T$ in standard coordinates,
while $(A^TA)^{-1}A^T$ is the pseudoinverse --- a different object.

\divider

%----------------------------------------------------------
\subsection*{Problem: Coordinates in an orthonormal basis and Parseval's identity}
%----------------------------------------------------------

Let $\mathcal{B} = \{\mathbf{q}_1, \mathbf{q}_2, \mathbf{q}_3\}$ be an orthonormal basis
for $\mathbb{R}^3$.
Suppose $\mathbf{v} \in \mathbb{R}^3$ satisfies
$$\langle \mathbf{v}, \mathbf{q}_1 \rangle = 4, \qquad
  \langle \mathbf{v}, \mathbf{q}_2 \rangle = -1, \qquad
  \langle \mathbf{v}, \mathbf{q}_3 \rangle = 3.$$
What is $\|\mathbf{v}\|^2$?

\begin{enumerate}[(A)]
\item \textcolor{blue}{\textbf{$26$}}
\item $16$
\item $6$
\item $36$
\item Cannot be determined without knowing the vectors $\mathbf{q}_i$ explicitly
\end{enumerate}

\textbf{Explanation:}
Since $\mathcal{B}$ is orthonormal, the coordinates of $\mathbf{v}$ \emph{are} the
inner products: $[\mathbf{v}]_\mathcal{B} = (4,-1,3)^T$.
By Parseval's identity,
$\|\mathbf{v}\|^2 = 4^2 + (-1)^2 + 3^2 = 16 + 1 + 9 = 26$.

\textbf{Trick answer:}
The answer $16$ squares only the largest coordinate (or stops after the first
term).

\textbf{Trick answer:}
The answer $6$ adds the coordinates ($4+(-1)+3=6$), confusing $\|\mathbf{v}\|^2$
with the sum of the coordinates.

\textbf{Partial credit:}
The answer $36$ computes $(4+(-1)+3)^2 = 36$ --- ``squares the sum instead of
summing the squares.''
The correct procedure is understood; the arithmetic rule is not.

\textbf{Trick answer:}
Saying the answer cannot be determined without knowing the $\mathbf{q}_i$ explicitly
misses the key advantage of an orthonormal basis: the norm can be read off
from the coordinates via Parseval's identity, regardless of which specific
orthonormal basis is used.

\divider

%%%%%%%%%%%%%%%%%%%%%%%%%%%%%%%%%%%%%%%%%%%%%%%%%%%%%%%%%%%
\section*{WEEKS 5/6: GRAM--SCHMIDT AND QR}
%%%%%%%%%%%%%%%%%%%%%%%%%%%%%%%%%%%%%%%%%%%%%%%%%%%%%%%%%%%

%----------------------------------------------------------
\subsection*{Problem: Gram--Schmidt --- what is guaranteed}
%----------------------------------------------------------

Let $V$ be an $n$-dimensional inner product space, and let
$S = \{\mathbf{v}_1, \ldots, \mathbf{v}_k\}$ be a linearly independent set in $V$
with $1 < k < n$.
The vectors in $S$ are \emph{not} pairwise orthogonal.
Which of the following is guaranteed?

\begin{enumerate}[(A)]
\item \textcolor{blue}{\textbf{There exists an orthonormal set of $k$ vectors whose span equals $\operatorname{span}(S)$.}}
\item There exist positive scalars $\alpha_i$ such that $\{\alpha_1\mathbf{v}_1, \ldots, \alpha_k\mathbf{v}_k\}$ is orthonormal.
\item There exists an orthonormal basis for $V$ that contains every vector in $S$ as an element.
\item The set $S$ can be orthonormalized only after first extending it to a basis for all of $V$.
\item None of the above is guaranteed.
\end{enumerate}

\textbf{Explanation:}
Gram--Schmidt applied to $\mathbf{v}_1,\ldots,\mathbf{v}_k$ produces orthonormal vectors
$\mathbf{q}_1,\ldots,\mathbf{q}_k$ satisfying
$\operatorname{span}\{\mathbf{q}_1,\ldots,\mathbf{q}_j\} = \operatorname{span}\{\mathbf{v}_1,\ldots,\mathbf{v}_j\}$
for each $j$.
In particular, $\operatorname{span}\{\mathbf{q}_1,\ldots,\mathbf{q}_k\} = \operatorname{span}(S)$.
No extension to a full basis is needed, and the original vectors need not
already be orthogonal.

\textbf{Trick answer:}
Rescaling each $\mathbf{v}_i$ by a positive scalar $\alpha_i$ preserves its
direction.
Normalization makes each vector unit length but \emph{cannot change angles}
between vectors.
Since the $\mathbf{v}_i$ are explicitly not pairwise orthogonal, no rescaling can
make them orthonormal.
This catches the ``Gram--Schmidt is just normalization'' misconception.

\textbf{Trick answer:}
Since at least two of the $\mathbf{v}_i$ satisfy $\langle\mathbf{v}_i,\mathbf{v}_j\rangle\neq 0$,
they cannot both belong to \emph{any} orthonormal set.
The student who picks this choice remembers the ``independent vectors extend
to a basis'' theorem but forgets that orthonormal extension requires modifying
the original vectors.

\textbf{Trick answer:}
Gram--Schmidt works on any linearly independent set of size $k \leq n$;
it orthonormalizes the span of $S$, not the full space.
No prior extension to $n$ vectors is needed.

\divider

%----------------------------------------------------------
\subsection*{Problem: Gram--Schmidt and the QR decomposition}
%----------------------------------------------------------

Let $A$ be an $n\times n$ invertible matrix with columns
$\mathbf{a}_1,\ldots,\mathbf{a}_n$, and let $A = QR$ be its QR decomposition.
Which statement best describes the relationship between Gram--Schmidt and QR?

\begin{enumerate}[(A)]
\item The columns of $Q$ are obtained by normalizing the columns of $A$; $R$ records the original norms.
\item \textcolor{blue}{\textbf{The columns of $Q$ are the Gram--Schmidt orthonormalization of $\mathbf{a}_1, \ldots, \mathbf{a}_n$; the entry $r_{ij}$ of $R$ is $\langle \mathbf{a}_j, \mathbf{q}_i \rangle$ for $i \leq j$.}}
\item The columns of $R$ are the Gram--Schmidt orthonormalization of $\mathbf{a}_1, \ldots, \mathbf{a}_n$; $Q$ records the change of basis.
\item Gram--Schmidt produces $Q$; the matrix $R = Q^T A$ is diagonal with the norms $\|\mathbf{a}_i\|$ on the diagonal.
\item Gram--Schmidt and QR solve different problems: Gram--Schmidt orthogonalizes vectors while QR factors a matrix.
\end{enumerate}

\textbf{Explanation:}
Gram--Schmidt applied to $\mathbf{a}_1,\ldots,\mathbf{a}_n$ produces orthonormal vectors
$\mathbf{q}_1,\ldots,\mathbf{q}_n$, which become the columns of $Q$.
Since each $\mathbf{a}_j$ is a linear combination of $\mathbf{q}_1,\ldots,\mathbf{q}_j$
(only earlier vectors appear in the orthogonalization of $\mathbf{a}_j$), we have
$\mathbf{a}_j = \sum_{i=1}^{j} r_{ij}\mathbf{q}_i$ with $r_{ij} = \langle\mathbf{a}_j,\mathbf{q}_i\rangle$.
This is exactly $A = QR$ with $R$ upper triangular.

\textbf{Trick answer:}
Normalizing without orthogonalizing does not produce orthonormal columns; the
projection-subtraction step is the heart of Gram--Schmidt.

\textbf{Wrong:}
This choice swaps the roles of $Q$ and $R$.

\textbf{Partial credit:}
The formula $R = Q^T A$ is correct (since $Q^T Q = I$, so $Q^T A = R$), but $R$
is upper triangular, not diagonal.
$R$ would be diagonal only if the columns of $A$ were already orthogonal.

\textbf{Wrong:}
QR and Gram--Schmidt are the same procedure viewed differently; QR is the
matrix packaging of Gram--Schmidt.

\divider

%----------------------------------------------------------
\subsection*{Problem: Using QR to solve a linear system}
%----------------------------------------------------------

An $n\times n$ invertible matrix $A$ has QR decomposition $A = QR$.
Using this factorization, the system $A\mathbf{x} = \mathbf{b}$ reduces to:

\begin{enumerate}[(A)]
\item Solving the lower triangular system $R^T\mathbf{x} = Q\mathbf{b}$.
\item Solving the upper triangular system $R\mathbf{x} = Q\mathbf{b}$.
\item Computing $\mathbf{x} = R^T Q^T \mathbf{b}$.
\item Computing $\mathbf{x} = QR^{-1}\mathbf{b}$.
\item \textcolor{blue}{\textbf{Solving the upper triangular system $R\mathbf{x} = Q^T\mathbf{b}$.}}
\end{enumerate}

\textbf{Explanation:}
From $QR\mathbf{x} = \mathbf{b}$, multiply both sides on the left by $Q^T$.
Since $Q$ is orthogonal, $Q^T Q = I$, giving $R\mathbf{x} = Q^T\mathbf{b}$.
The right-hand side $Q^T\mathbf{b}$ is a single matrix-vector product (cheap),
and $R$ is upper triangular, so the resulting system is solved by
back-substitution.
This is the practical payoff of QR: it replaces a general linear system with a
triangular one.

\textbf{Partial credit:}
The choice solving the upper triangular system $R\mathbf{x} = Q\mathbf{b}$ correctly
identifies the structure (triangular system), but uses $Q$ instead of $Q^T$ on
the right-hand side.
This student understands the triangular structure but has not internalized
$Q^{-1} = Q^T$, the defining property of orthogonal matrices.

\textbf{Trick answer:}
The choice with $R^T$ (lower triangular) and $Q$ transposes the wrong factor.

\textbf{Trick answer:}
$R^T Q^T = (QR)^T = A^T$, so $R^T Q^T\mathbf{b} = A^T\mathbf{b}$, not $A^{-1}\mathbf{b}$.
This is the student who confuses transpose with inverse.

\textbf{Trick answer:}
$A^{-1} = R^{-1}Q^T$, not $QR^{-1}$.
Products invert in reverse order: $(QR)^{-1} = R^{-1}Q^{-1}$.
Moreover, explicitly computing $R^{-1}$ defeats the purpose of the factorization.

\divider

%%%%%%%%%%%%%%%%%%%%%%%%%%%%%%%%%%%%%%%%%%%%%%%%%%%%%%%%%%%
\section*{WEEK 6: ORTHOGONAL DECOMPOSITION AND THE PSEUDOINVERSE}
%%%%%%%%%%%%%%%%%%%%%%%%%%%%%%%%%%%%%%%%%%%%%%%%%%%%%%%%%%%

%----------------------------------------------------------
\subsection*{Problem: The adjoint --- domain, codomain, and injectivity}
%----------------------------------------------------------

Let $T: \mathbb{R}^3 \to \mathbb{R}^2$ be defined by
$T(x_1, x_2, x_3) = (x_1 + 2x_3,\; x_2 - x_3)$,
with the standard dot product on both spaces.
Which statement about the adjoint $T^*$ is true?

\begin{enumerate}[(A)]
\item $T^*: \mathbb{R}^2 \to \mathbb{R}^3$ is surjective.
\item $T^*: \mathbb{R}^3 \to \mathbb{R}^2$ and has the same kernel as $T$.
\item $\operatorname{im}(T^*) = \ker(T)$.
\item \textcolor{blue}{\textbf{$T^*: \mathbb{R}^2 \to \mathbb{R}^3$ is injective.}}
\item $T^*$ does not exist because $T$ is not invertible.
\end{enumerate}

\textbf{Explanation:}
The matrix of $T$ is $A = \begin{bmatrix}1&0&2\\0&1&-1\end{bmatrix}$, so $T^*$ is
represented by $A^T = \begin{bmatrix}1&0\\0&1\\2&-1\end{bmatrix}$, mapping
$\mathbb{R}^2\to\mathbb{R}^3$.
This $3\times 2$ matrix has rank $2$ (both columns are linearly independent),
so $T^*$ is injective.

\textbf{Trick answer:}
$T^*$ maps into $\mathbb{R}^3$ with a $2$-dimensional image, so it is \emph{not}
surjective (its image misses a $1$-dimensional piece of $\mathbb{R}^3$).

\textbf{Trick answer:}
The choice getting the domain and codomain of $T^*$ backwards ($T^*: \mathbb{R}^3\to\mathbb{R}^2$)
also claims $\ker(T^*) = \ker(T)$, which is doubly wrong:
$\ker(T)$ is $1$-dimensional in $\mathbb{R}^3$ while $\ker(T^*) = \{\mathbf{0}\}$.

\textbf{Partial credit:}
The choice $\operatorname{im}(T^*) = \ker(T)$ recalls the Fundamental Theorem of
Linear Algebra but states it incorrectly.
The correct relationship is $\operatorname{im}(T^*) = \ker(T)^\perp$.
Here $\ker(T)$ is $1$-dimensional while $\operatorname{im}(T^*)$ is $2$-dimensional;
they are orthogonal complements in $\mathbb{R}^3$, not equal.
This is the ``off by a complement'' error, showing genuine FTLA awareness with
an incomplete grasp of the orthogonal structure.

\textbf{Wrong:}
Adjoints always exist for linear maps between inner product spaces;
invertibility of $T$ is irrelevant.

\divider

%----------------------------------------------------------
\subsection*{Problem: The pseudoinverse --- geometric interpretation}
%----------------------------------------------------------

Let $T: V \to W$ be a linear transformation between inner product spaces
(not necessarily injective or surjective).
Given $\mathbf{b} \in W$, which description of $T^\dagger \mathbf{b}$ is correct?

\begin{enumerate}[(A)]
\item Project $\mathbf{b}$ onto $(\operatorname{im} T)^\perp$, then apply $T^{-1}$ to the result.
\item \textcolor{blue}{\textbf{Project $\mathbf{b}$ onto $\operatorname{im}(T)$; then map the projection back to the unique preimage in $(\ker T)^\perp$ via the invertible restriction $T\big|_{(\ker T)^\perp}$.}}
\item Apply the adjoint $T^*$ to $\mathbf{b}$; the adjoint always inverts $T$ on $\operatorname{im}(T)$.
\item Project $\mathbf{b}$ onto $\ker(T^*)$, then solve $T\mathbf{x} = \mathbf{b}$ in that subspace.
\item Find any solution to $T\mathbf{x} = \mathbf{b}$, then project $\mathbf{x}$ onto $\ker(T)$.
\end{enumerate}

\textbf{Explanation:}
The restriction of $T$ to $(\ker T)^\perp$ is an isomorphism onto
$\operatorname{im}(T)$.
The pseudoinverse exploits this in two steps:
(1) project $\mathbf{b}$ onto $\operatorname{im}(T)$, discarding the component in
$(\operatorname{im} T)^\perp = \ker(T^*)$;
(2) apply the inverse of the restricted map $T|_{(\ker T)^\perp}$ to pull the
projected vector back to the unique preimage in $(\ker T)^\perp$.
This yields the minimum-norm least-squares solution: ``minimum norm'' because
the output lives in $(\ker T)^\perp$, and ``least squares'' because the input
is the closest point in $\operatorname{im}(T)$ to $\mathbf{b}$.

\textbf{Trick answer:}
Projecting onto $(\operatorname{im} T)^\perp$ gives the part of $\mathbf{b}$ that $T$
\emph{cannot reach}, so there is no preimage to find there.
Also, $T^{-1}$ need not exist.

\textbf{Trick answer:}
$T^*$ maps $W\to V$ and sends $\operatorname{im}(T)$ to $(\ker T)^\perp$, but it does
not invert $T$; it distorts magnitudes.
Students who half-remember the role of $A^T$ in $A^\dagger = (A^TA)^{-1}A^T$
and attribute the inversion entirely to the adjoint pick this.

\textbf{Trick answer:}
$\ker(T^*) = (\operatorname{im} T)^\perp$ is exactly the component of $\mathbf{b}$ with
\emph{no} preimage under $T$ --- the opposite of what the pseudoinverse needs.

\textbf{Trick answer:}
Projecting a solution onto $\ker(T)$ gives the component that $T$ sends to
zero --- the opposite of the minimum-norm idea.
The pseudoinverse produces a vector in $(\ker T)^\perp$, not $\ker(T)$.

\divider

%----------------------------------------------------------
\subsection*{Problem: Orthogonal projection matrices --- which property is FALSE}
%----------------------------------------------------------

Let $P$ be the matrix that orthogonally projects $\mathbb{R}^n$ onto a subspace
$W$ with $\dim(W) = k$, where $0 < k < n$.
Which of the following is \emph{false}?

\begin{enumerate}[(A)]
\item $P^2 = P$
\item $P^T = P$
\item $\operatorname{rank}(P) = k$
\item $I - P$ is the orthogonal projection onto $W^\perp$
\item \textcolor{blue}{\textbf{$P$ is invertible}}
\end{enumerate}

\textbf{Explanation:}
Since $0 < k < n$, the projection has a nontrivial kernel ($W^\perp \neq \{\mathbf{0}\}$):
every $\mathbf{v} \in W^\perp$ satisfies $P\mathbf{v} = \mathbf{0}$.
Therefore $P$ is not invertible.
All other listed properties are fundamental facts about orthogonal projections.

\textbf{Wrong:}
$P^2 = P$ is the defining algebraic property of any projection: projecting twice
is the same as projecting once.

\textbf{Wrong:}
$P^T = P$ is what makes the projection \emph{orthogonal}; it distinguishes
orthogonal projections from oblique ones.

\textbf{Wrong:}
$\operatorname{im}(P) = W$, so $\operatorname{rank}(P) = \dim(W) = k$.

\textbf{Wrong:}
In the orthogonal decomposition $\mathbf{v} = P\mathbf{v} + (I-P)\mathbf{v}$, the term
$(I-P)\mathbf{v}$ lies in $W^\perp$, so $I-P$ is the projection onto $W^\perp$.

\divider

%----------------------------------------------------------
\subsection*{Problem: The four fundamental subspaces (FTLA)}
%----------------------------------------------------------

Let $A$ be a $5\times 8$ matrix with $\operatorname{rank}(A) = 3$.
Consider the four fundamental subspaces of $A$:
$\ker(A)$, $\operatorname{im}(A)$, $\ker(A)^\perp$, and $\operatorname{im}(A)^\perp$.
Which of the following is true?

\begin{enumerate}[(A)]
\item $\ker(A)$ and $\operatorname{im}(A)$ are orthogonal complements.
\item $\ker(A)$ has dimension $3$ and is a subspace of $\mathbb{R}^5$.
\item $\dim(\operatorname{im}(A)^\perp) = 3$
\item $\dim(\ker(A)^\perp) = 5$
\item \textcolor{blue}{\textbf{A vector $\mathbf{x} \in \mathbb{R}^8$ that is orthogonal to every row of $A$ must be in $\ker(A)$.}}
\end{enumerate}

\textbf{Explanation:}
The Fundamental Theorem of Linear Algebra (geometric form) states that
$\ker(A) = (\text{row space of }A)^\perp$ in $\mathbb{R}^8$.
A vector orthogonal to every row of $A$ lies in the orthogonal complement of
the row space, which is exactly $\ker(A)$.
This is the orthogonal restatement of the fact that $A\mathbf{x} = \mathbf{0}$ if and
only if each row of $A$ dots to zero with $\mathbf{x}$.

\textbf{Trick answer:}
$\ker(A) \subset \mathbb{R}^8$ and $\operatorname{im}(A) \subset \mathbb{R}^5$ live in
\emph{different ambient spaces}, so they cannot be orthogonal complements.
The correct complementary pairs are:
$\ker(A) \oplus \ker(A)^\perp = \mathbb{R}^8$ and
$\operatorname{im}(A) \oplus \operatorname{im}(A)^\perp = \mathbb{R}^5$.

\textbf{Wrong:}
Two errors: $\dim(\ker(A)) = 8 - 3 = 5$, not $3$; and $\ker(A) \subset \mathbb{R}^8$,
not $\mathbb{R}^5$.

\textbf{Trick answer:}
$\operatorname{im}(A) \subset \mathbb{R}^5$ has dimension $3$, so
$\dim(\operatorname{im}(A)^\perp) = 5-3 = 2$, not $3$.
Students who assume a subspace and its complement have equal dimension pick this.

\textbf{Trick answer:}
$\ker(A) \subset \mathbb{R}^8$ has dimension $5$, so
$\dim(\ker(A)^\perp) = 8 - 5 = 3$, not $5$.
Students who confuse the dimension of the subspace with the dimension of its
complement pick this.

\divider

%----------------------------------------------------------
\subsection*{Problem: The pseudoinverse as a left inverse}
%----------------------------------------------------------

Let $A$ be an $m\times n$ matrix with $m > n$ and $\operatorname{rank}(A) = n$
(full column rank).
The pseudoinverse is $A^\dagger = (A^TA)^{-1}A^T$.
Which of the following statements is TRUE?

\begin{enumerate}[(A)]
\item $A^\dagger A = I_m$, so $A^\dagger$ is a right inverse of $A$.
\item \textcolor{blue}{\textbf{$A^\dagger$ is an $n\times m$ matrix, and $A^\dagger A = I_n$.}}
\item $(A^TA)^{-1}$ does not exist because $A$ is not square.
\item $A^\dagger \mathbf{b}$ solves $A\mathbf{x} = \mathbf{b}$ exactly for every $\mathbf{b} \in \mathbb{R}^m$.
\item $A^\dagger = A^{-1}$ after discarding the extra rows of $A$.
\end{enumerate}

\textbf{Explanation:}
Since $A$ is $m\times n$, $A^T$ is $n\times m$, and $A^TA$ is $n\times n$.
Full column rank of $A$ guarantees $A^TA$ is invertible, so $(A^TA)^{-1}$
exists.
Thus $A^\dagger = (A^TA)^{-1}A^T$ is $n\times m$.
Computing $A^\dagger A = (A^TA)^{-1}(A^TA) = I_n$: the pseudoinverse is a
\emph{left} inverse of $A$.

\textbf{Partial credit:}
The choice claiming $A^\dagger A = I_m$ correctly recognizes that $A^\dagger$
acts as a one-sided inverse yielding an identity matrix, but confuses left with
right.
Note also the dimension mismatch: $A^\dagger A$ is $n\times n$, so the result
cannot be $I_m$ (which is $m\times m$).

\textbf{Trick answer:}
$A^TA$ is $n\times n$ (not $m\times m$), and full column rank of $A$ guarantees
it is invertible.
Students who worry that $A^TA$ might be defective because $A$ is not square
have not tracked the dimensions correctly.

\textbf{Trick answer:}
$A\mathbf{x} = \mathbf{b}$ is overdetermined ($m > n$), so exact solutions exist only
when $\mathbf{b}\in\operatorname{im}(A)$.
For general $\mathbf{b}$, $A^\dagger\mathbf{b}$ gives the least-squares solution
minimizing $\|A\mathbf{x}-\mathbf{b}\|$, not an exact solution.

\textbf{Wrong:}
There is no valid operation of ``discarding the extra rows to produce an
inverse''; the pseudoinverse accounts for the rectangular shape through the
Gram matrix $A^TA$, not by truncation.

\divider

%----------------------------------------------------------
\subsection*{Problem: Orthogonal projection and the residual}
%----------------------------------------------------------

Let $W$ be a subspace of an inner product space $V$, and let $\mathbf{v}\in V$.
Define $\hat{\mathbf{v}} = \Pi_W(\mathbf{v})$ (the orthogonal projection onto $W$) and
$\mathbf{r} = \mathbf{v} - \hat{\mathbf{v}}$ (the ``error'').
Which of the following is \emph{guaranteed} to be true?

\begin{enumerate}[(A)]
\item $\|\hat{\mathbf{v}}\| = \|\mathbf{v}\|$
\item $\langle \hat{\mathbf{v}},\, \mathbf{r} \rangle = 1$
\item $\mathbf{r} \in W$
\item $\|\mathbf{r}\| > 0$
\item \textcolor{blue}{\textbf{None of the above}}
\end{enumerate}

\textbf{Explanation:}
The orthogonal decomposition gives $\mathbf{v} = \hat{\mathbf{v}} + \mathbf{r}$ with
$\hat{\mathbf{v}} \in W$ and $\mathbf{r} \in W^\perp$, satisfying
$\langle\hat{\mathbf{v}},\mathbf{r}\rangle = 0$.
By the Pythagorean theorem, $\|\mathbf{v}\|^2 = \|\hat{\mathbf{v}}\|^2 + \|\mathbf{r}\|^2$,
so $\|\hat{\mathbf{v}}\| \leq \|\mathbf{v}\|$ with equality if and only if
$\mathbf{v}\in W$.
None of (A)--(D) is universally true.

\textbf{Trick answer:}
$\|\hat{\mathbf{v}}\| = \|\mathbf{v}\|$ holds only when $\mathbf{v}\in W$ (so $\mathbf{r} = \mathbf{0}$).
For any $\mathbf{v}\notin W$, Pythagoras gives strict inequality.

\textbf{Trick answer:}
$\langle\hat{\mathbf{v}},\mathbf{r}\rangle = 0$, not $1$: orthogonality is the
defining geometric property of the decomposition.

\textbf{Partial credit:}
$\mathbf{r}$ belongs to $W^\perp$, not $W$.
This is the most common swap error in orthogonal decomposition.
Choosing the answer that says $\mathbf{r}\in W$ shows awareness that $\mathbf{r}$ lives
in a specific subspace related to $W$, but gets the wrong one.

\textbf{Trick answer:}
$\|\mathbf{r}\| > 0$ holds for ``generic'' $\mathbf{v}$ and in the overdetermined
least-squares setting students encounter most often, but fails when $\mathbf{v}\in W$:
then $\hat{\mathbf{v}} = \mathbf{v}$ and $\mathbf{r} = \mathbf{0}$.

\end{document}